% !TEX program = pdflatex
\documentclass[12pt,a4paper]{letter}

\usepackage[T1]{fontenc}
\usepackage[margin=1in]{geometry}
\usepackage{hyperref}
\usepackage{setspace}
\usepackage{parskip}

\hypersetup{
  colorlinks=true,
  urlcolor=blue
}

\onehalfspacing

\signature{Professor Ayo Ogunjuyigbe\\
Professor of Electrical \& Electronic Engineering\\
Provost, Postgraduate School\\
University of Ibadan, Nigeria}

\address{Postgraduate School\\
University of Ibadan\\
Ibadan, Oyo State\\
Nigeria}

\date{27 February 2026}

\begin{document}

\begin{letter}{UK Visas \& Immigration\\
Global Talent Visa --- Tech Nation (Digital Technology)\\
United Kingdom}

\opening{To Whom It May Concern,}

I write to endorse Odeyemi Olusegun Israel for the UK Global Talent Visa
under the Digital Technology route.  I am a Professor of Electrical and
Electronic Engineering at the University of Ibadan, where I currently
serve as Provost of the Postgraduate School.  My research spans power
systems, computational intelligence, and applied machine learning.  I
have mentored Olusegun for over 10~years and have reviewed his technical
work in detail.

%% ── AMTTP ──────────────────────────────────────────────────────────

\textbf{AMTTP: a production compliance platform built by one person.}

Olusegun has designed and built the Anti-Money Laundering Transaction
Trust Protocol (AMTTP)---a four-layer system comprising client SDKs
(TypeScript and Python, both open-sourced), a compliance orchestration
engine of 7~parallel microservices, an ML training pipeline processing
2.64~million Ethereum transactions (production ROC-AUC of~0.9879), and
18~Solidity smart contracts on Ethereum Sepolia.  The whole stack runs
across 17~containerised microservices.

AMTTP solves a problem the blockchain compliance industry has not yet
solved: \textit{pre-settlement} AML enforcement.  Every major tool on
the market---Chainalysis, Elliptic, TRM~Labs---operates after a
transaction has settled.  AMTTP intervenes before settlement, completing
ML inference, graph analysis, sanctions screening, and policy evaluation
within Ethereum's 12-second block window.  That is genuine real-time
systems engineering.

He also built zkNAF, a zero-knowledge proof framework (Groth16
ZK-SNARKs on BN254) that lets institutions prove compliance
facts---KYC validity, risk-score range, sanctions
non-membership---without revealing personal data, resolving the
GDPR-versus-AML conflict at the cryptographic level.

In a university, a project of this scope would require a PI, several
postdocs, and a team of PhD students.  Olusegun built it alone.  I have
reviewed the GitHub repository: there is no evidence of collaborators or
outsourced components.

%% ── SIAM PAPER / ERCOT ────────────────────────────────────────────

\textbf{Blind-Spot Decomposition and the Geometry of System Collapse: a
contribution I can assess directly.}

Olusegun's research papers on Blind Spot Decomposition Theory (BSDT) and
the Universal Detection Principle (UDP) are solid contributions to
anomaly detection.  But it is his third paper, \textit{Blind-Spot
Decomposition and the Geometry of System Collapse}, that I endorse most
directly, because it contains results in power systems that fall
squarely within my area of expertise.

In February~2021, Winter Storm Uri caused cascading blackouts across the
Texas ERCOT grid.  The public narrative blamed frozen wind turbines.
Olusegun's adaptive friction framework, applied to ERCOT dispatch data,
achieves:

\begin{itemize}
  \item \textbf{Detection 72~hours before rolling blackouts began}
    (AUC~=~0.9997).
  \item \textbf{Adaptive friction reduces peak capacity loss by 71\%.}
  \item Per-agent BSDT decomposition identifies \textbf{gas supply
    chain disruption}---not frozen wind turbines---as the structural
    vulnerability driving the cascade.
\end{itemize}

The gas supply finding is exactly what the subsequent FERC/NERC joint
investigation confirmed---but Olusegun's framework identifies it
\textit{mathematically, 72~hours in advance, from dispatch data alone}.
As a power systems researcher, I can say that a 72-hour predictive lead
on a grid cascade event, with correctly identified root cause, is an
exceptional result.  I am not aware of any other framework in the
literature that achieves this.

The same paper validates on FDIC banking data (6-quarter lead on the
2008~GFC, zero false alarms), the Terra/Luna collapse ($r = 0.996$
correlation), and 3D Navier-Stokes fluid dynamics.  That the same
mathematical framework diagnoses blockchain fraud, predicts banking
crises, and identifies power grid cascade origins is remarkable.

%% ── ASSESSMENT ─────────────────────────────────────────────────────

\textbf{Relevance to the UK and my assessment.}

The FCA's forthcoming crypto-asset regulatory framework will create
demand for protocol-level compliance infrastructure.  AMTTP is, as far
as I am aware, the only system providing pre-settlement AML enforcement
on blockchain transactions.  The ERCOT and Navier-Stokes work also
aligns with UK priorities in infrastructure resilience (Ofgem, National
Grid ESO) and applied mathematics (EPSRC).

In my years at the University of Ibadan, I have mentored hundreds of
engineers.  Olusegun's engineering capability and systems-level thinking
place him in the top tier of talent I have encountered.  He designs and
builds complex systems from first principles---across ML, blockchain,
cryptography, and mathematical research---entirely without institutional
support.

I recommend him for the Global Talent Visa without reservation.

\closing{Yours faithfully,}

\end{letter}

\end{document}
